\documentclass[11pt]{article}
\usepackage[T1]{fontenc}
\usepackage{textcomp}
\usepackage[margin=0.75in]{geometry}
\usepackage{amsmath,amssymb,amsthm}
\usepackage{array,booktabs}
\usepackage{hyperref}
\setlength{\parindent}{0pt}
\newtheorem{theorem}{Theorem}
\theoremstyle{definition}
\newtheorem{definition}{Definition}
\title{Flow Proof Artifact: Nat-mul-zero-left}
\author{Generated by \texttt{flow doc proof}}
\date{Flow Proof Book}
\begin{document}
\maketitle
Zero on the left annihilates multiplication.\\[0.5em]
\textbf{Source.} peano — https://en.wikipedia.org/wiki/Peano\_axioms\\[0.5em]
\bigskip
\noindent{\large \textbf{Derived fact 1.} \textit{0 * n = 0} \hfill the natural numbers \cdot multiplication \cdot zero on the left gives zero}\par
\smallskip\noindent\textit{Coordinate: the natural numbers \cdot multiplication \cdot zero on the left gives zero}\par
\smallskip\noindent\textit{Source: peano}\par
\smallskip\noindent\textit{Built on: zero is the left annihilator, for multiplication on the natural numbers, zero is the right annihilator, for multiplication on the natural numbers, order does not matter, for multiplication on the natural numbers}\par
\medskip
\noindent\textbf{Goal.} 0 * n = 0\par
\noindent$\displaystyle \forall n \in \mathbb{N}\quad 0 \cdot n = 0$\par
\medskip
\noindent
\begin{tabular}{@{} >{\raggedright\arraybackslash}p{0.06\textwidth} >{\raggedright\arraybackslash}p{0.47\textwidth} >{\raggedleft\arraybackslash}p{0.41\textwidth} @{}}
\textbf{\#} & \textbf{Proof} & \textbf{Mathematics} \\
\midrule
\textbf{1.} & We prove that zero on the left gives zero for multiplication on the natural numbers. & \\[0.45em]
\textbf{2.} & We invoke the definitional clause governing multiplication on the natural numbers: zero is the left annihilator, for multiplication on the natural numbers (instantiated for n). & \\[0.45em]
\textbf{3.} & We invoke the definitional clause governing multiplication on the natural numbers: zero is the right annihilator, for multiplication on the natural numbers (instantiated for 0). & \\[0.45em]
\textbf{4.} & We invoke the derived fact governing multiplication on the natural numbers: order does not matter, for multiplication on the natural numbers (instantiated for 0, n). & \\[0.45em]
\textbf{5.} & From step 2, step 3, and step 4, this implies 0 times n equals 0. Hence proven. & $\displaystyle 0 \cdot n = 0$ \\[0.45em]
\bottomrule
\end{tabular}
\label{thm:Nat-multiplication-zero_on_the_left_gives_zero}
\par\medskip\hrule
\end{document}
